% from ega1-1-fr.tex
\oldpage[I]{92}
prolonge en une section $s$ de $\mathscr{F}$ sur $V$~; appliquant
\textit{d}\,2), on voit qu'il existe un entier $n\geq0$ tel que
$f^n g^m t=0$, et comme l'image de $g$ dans $A_g$ est inversible, $f^n t=0$.

Cela étant, comme $V$ est quasi-compact, le lemme
(\hyperref[I.1.4.1.1-fr]{1.4.1.1}) prouve que les conditions
\textit{d}'\,1) et \textit{d}'\,2) sont vérifiées. Considérons alors le
$A$-module $M=\Gamma(V,\mathscr{F})$, et définissons un homomorphisme de
$\widetilde{A}$-Modules $u:\widetilde{M}\to j_*(\mathscr{F})$, où $j$ est
l'injection canonique $V\to X$. Comme les $D(f)$ forment une base de la
topologie de $X$, il suffit, pour chaque $f\in A$, de définir un
homomorphisme
$u_f:M_f\to\Gamma(D(f),j_*(\mathscr{F}))=\Gamma(D(f)\cap V,\mathscr{F})$,
avec les conditions de compatibilité usuelles
(\hyperref[0.3.2.5-fr]{0, 3.2.5}). Comme l'image canonique de $f$ dans
$A_f$ est inversible, l'homomorphisme de restriction
$M=\Gamma(V,\mathscr{F})\to\Gamma(D(f)\cap V,\mathscr{F})$ se factorise en
$M\longrightarrow M_f\xrightarrow{u_f}\Gamma(D(f)\cap V,\mathscr{F})$
(\hyperref[0.1.2.4-fr]{0, 1.2.4}), et la vérification des conditions de
compatibilité pour $D(g)\subset D(f)$ est immédiate. Cela étant, montrons
que la condition \textit{d}'\,1) (resp. \textit{d}'\,2)) entraîne que chacun
des $u_f$ est surjectif (resp. injectif), ce qui prouvera que $u$ est
\emph{bijectif}, et par suite que $\mathscr{F}$ est la restriction à $V$
d'un $\widetilde{A}$-Module isomorphe à $\widetilde{M}$. Or, si
$s\in\Gamma(D(f)\cap V,\mathscr{F})$, il existe d'après \textit{d}'\,1) un
entier $n\geq0$ tel que $f^n s$ se prolonge en une section $z\in M$~; on a
alors $u_f(z/f^n)=s$, donc $u_f$ est surjectif. De même, si $z\in M$ est tel
que $u_f(z/1)=0$, cela signifie que la restriction à $D(f)\cap V$ de la
section $z$ est nulle~; d'après \textit{d}'\,2), il existe un entier
$n\geq0$ tel que $f^n z=0$, d'où $z/1=0$ dans $M_f$, et $u_f$ est donc
injectif.

\par\hfill C.Q.F.D.\par

\begin{corollary}[1.4.2]
\label{I.1.4.2-fr}
Tout faisceau quasi-cohérent sur un ouvert quasi-compact de $X$ est induit
par un faisceau quasi-cohérent sur $X$.
\end{corollary}

\begin{corollary}[1.4.3]
\label{I.1.4.3-fr}
Toute $\mathscr{O}_X$-Algèbre quasi-cohérente sur
$X=\operatorname{Spec}(A)$ est isomorphe à une $\mathscr{O}_X$-Algèbre de
la forme $\widetilde{B}$, où $B$ est une algèbre sur $A$~; tout
$\widetilde{B}$-Module quasi-cohérent est isomorphe à un
$\widetilde{B}$-Module de la forme $\widetilde{N}$, où $N$ est un
$B$-module.
\end{corollary}

En effet, une $\mathscr{O}_X$-Algèbre quasi-cohérente est un
$\mathscr{O}_X$-Module quasi-cohérent, donc de la forme $\widetilde{B}$, où
$B$ est un $A$-module~; le fait que $B$ soit une $A$-algèbre résulte de la
caractérisation de la structure de $\mathscr{O}_X$-Algèbre à l'aide de
l'homomorphisme
$\widetilde{B}\otimes_{\widetilde{A}}\widetilde{B}\to\widetilde{B}$ de
$\widetilde{A}$-Modules, ainsi que de
(\hyperref[I.1.3.12-fr]{1.3.12}). Si $\mathscr{G}$ est un
$\widetilde{B}$-Module quasi-cohérent, il suffit de montrer que $\mathscr{G}$
est aussi un $\widetilde{A}$-Module quasi-cohérent pour conclure ensuite de
la même façon~; comme la question est locale, on peut, en se restreignant à
un ouvert de $X$ de la forme $D(f)$, supposer que $\mathscr{G}$ est le
conoyau d'un homomorphisme
$\widetilde{B}^{(I)}\to\widetilde{B}^{(J)}$ de $\widetilde{B}$-Modules (et
\emph{a fortiori} de $\widetilde{A}$-Modules)~; la proposition résulte alors
de (\hyperref[I.1.3.8-fr]{1.3.8}) et
(\hyperref[I.1.3.9-fr]{1.3.9}).

\subsection{Faisceaux cohérents sur un spectre premier.}
\label{subsection:I.1.5-fr}

\begin{theorem}[1.5.1]
\label{I.1.5.1-fr}
Soient $A$ un anneau \emph{noethérien}, $X=\operatorname{Spec}(A)$ son spectre
premier, $V$ une partie ouverte de $X$, $\mathscr{F}$ un
$(\mathscr{O}_X|V)$-Module. Les conditions suivantes sont équivalentes~:
\begin{enumerate}
  \item[\textit{a})] $\mathscr{F}$ est cohérent.
  \item[\textit{b})] $\mathscr{F}$ est de type fini et quasi-cohérent.
  \item[\textit{c})] Il existe un $A$-module $M$ de type fini tel que
  $\mathscr{F}$ soit isomorphe au faisceau $\widetilde{M}|V$.
\end{enumerate}
\end{theorem}

